\beginsong{Prière scoute}[by={Saint Ignace de Loyola, Gaston Schindler}]
\beginverse
\[A]Seigneur Jé\[Bm]sus,
\[A]Appre\[D]nez-nous à \[E7]être gé\[A]néreux,
À vous ser\[D]vir comme \[E7]vous le mé\[A]ri\[C#]tez,
\[F#m]À don\[C#m]ner sans \[D]comp\[F#7]ter,
\endverse
\beginverse
\[Bm]À com\[Bm7]battre sans \[E7]souci des bles\[C#7]su\[F#m]res,
À tra\[B7]vail\[Em]ler sans cher\[A]cher le re\[F#m]pos,
\[D]À nous dé\[C#m]pen\[D]ser sans at\[C#7]tendre d'au\[D]tre ré\[F#7]com\[Bm]pense que
\[F#m]cel\[E7]le de sa\[A]voir que nous fai\[F#m]sons
\[E7]vo\[D]tre sain\[E]te vo\[A]lonté.
\endverse
\endsong

\beginsong{Viens grandir avec les farfadets}[by={Romain Avice, Sylvain Bréant}]
\textnote{Capo 2}
\beginverse
\[C]Dans la nature, je \[Am]vis de belles aventures
\[F]J'éveille mes sens, \[C]expé\[G]rience
\[C]Je cours, je ris, avec \[Am]mes bras je construis
\[F]Fille ou garçon, \[C]cré\[G]ation
\[C]Rando fatigante, j'en\[Am]tends les oiseaux qui chantent,
\[F]Que c'est beau la vie, \[C]je sou\[G]ris !
\[C]Questions en tête, en fai\[Am]sant des pirouettes,
\[F]Pas d'Internet, \[C]c'est la \[G]fête
\endverse
\beginverse
\[Dm]Tom et Prune sont nos a\[G]mis
\[Dm]En vert pomme, réu\[G]nis
\endverse
\beginchorus
\[C]Viens grandir avec les \[G]Farfadets
\[Am]L'épopée d'la ronde est \[Em]en projet
\[F]Les parents plantent \[G]le décor
\[F]Tom et Prune sortent le \[C]tré\[G]sor, hey !
\[C]Viens grandir avec les \[G]Farfadets
\[Am]N'oublie pas d'apporter \[Em]ton carnet
\[F]Réunis li\[G]sons la loi,
\[F]Chantons tous d'une \[C]même \[G]voix, tralalala
\endchorus
\beginverse
\[C]J'ouvre ma porte à ce que \[Am]l'autre m'apporte
\[F]Les différences, \[C]c'est une \[G]chance
\[C]Ensemble on joue, et on \[Am]s'amuse comme des fous
\[F]Bientôt on d'vient, \[C]des co\[G]pains
\[C]Donne ton avis, et c'est \[Am]comme ça qu'on choisit
\[F]Dans le respect, \[C]c'est la \[G]paix
\[C]Servir on assure, et on \[Am]protège la nature,
\[F]Sortir d'sa case, \[C]c'est la \[G]BASE
\endverse
\beginverse
\[C]Sens dans ton cœur, un arc \[Am]en ciel de couleurs
\[F]Identifions, \[C]émo\[G]tions
\[C]Chante à Jésus, tout ce \[Am]que l'on a vécu
\[F]L'amour rayonne, \[C]je par\[G]donne
\[C]Je suis comme je suis, dans \[Am]ma ronde d'amis
\[F]Jésus merci, \[C]accueil\[G]li
\[C]Que c'est épatant, de dé\[Am]couvrir ses talents
\[F]Viens dans la ronde, \[C]change le \[G]monde
\endverse
\endsong

\beginsong{La ronde des Farfadets}[by={Raphaël Michel}]
\beginverse
\[Bb]Avec tous mes a\[Bb]mis
Tout au \[Bb]long de la \[Bb]vie
Les \[Gm]far-fa-\[Cm]dets sont \[Eb]là
Pour me don\[Bb]ner la \[F]joie
\endverse
\beginverse
\[Gm]Ils m'aident à mieux \[Cm]co\[Eb]nnaître
Mes mains mon \[Bb]cœur ma \[F]tê\[Bb]te
\[Eb]Nous chantons à tue-\[Bb]tête
Dans les jeux et la \[F7]fête
\endverse
\beginchorus
Nous voi\[Bb]ci dans la \[Eb]ronde nous sommes Farfa\[Bb]det
Pour a\[F7]gir dans le mon\[Bb]de tou\[Eb]jours en se\[Bb]cret
Nous voi\[F7]ci dans la \[Bb]ronde nous sommes Farfa\[Eb]det
Nous ai\[Bb]dons tout le \[F7]mon\[Bb]de
\endchorus
\beginverse
\[Bb]Quand quelqu'un près de moi
A besoin d'un peu d'aide
Quand quelqu'un près de moi
A plus besoin qu'on l'aime
\endverse
\beginverse
\[Bb]Ce sont les farfadets
Qui en secret m'apprennent
Qu'on ne laisse jamais
Un ami dans la peine\ldots
\endverse
\beginverse
\[Bb]Avec tous mes amis
Ils m'aident à partager
A embellir ma vie
Pour toujours mieux aimer
\endverse
\beginverse
\[Bb]C'est grâce aux farfadets
Que j'aime découvrir
Tout ce que Dieu a fait
La nature et les rires
\endverse
\endsong

\beginsong{Farfadets, un nouveau monde}[by={Jean Jourdier}]
\textnote{Capo 2}
\beginverse
\[C]Connaissez-vous l'his\[G]toire
De \[Am]petits lutins \[G]verts
Qui \[Am]voulaient vraiment y \[G]croire
Aux \[C]hommes sur la \[G]terre
\[C]Ils rêvaient de grands \[G]voyages
Et de \[Am]bateaux sur la \[G]mer
De \[Am]voir de beaux pay\[G]sages
Bri\[C]ser les barri\[G]ères
\endverse
\beginverse*
\[Dm]Si vous voulez savoir qui ils \[G]sont
Et \[Am]connaître leur uni\[G]vers
Il \[Am]suffit de crier leur \[G]nom
Et \[Am]lever les bras en l'\[G]air
\endverse*
\beginchorus
\[C]Farfadets, un nouveau \[G]monde
\[Am]Je profite de chaque se\[G]conde
\[Am]Je t'entraîne dans ma \[G]ronde
\[C]Où je te prends la \[G]main
Iheo, iheo, iheo, iheee !
\endchorus
\beginverse
\[C]Ils grandissent dans la fo\[G]rêt,
Leur ami\[Am]tié vaut de l'\[G]or.
Si ja\[Am]mais vous décou\[G]vrez
Leur \[C]Malle aux tré\[G]sors,
\[C]Vous y trouverez des \[G]rêves,
Des re\[Am]gards et des sou\[G]rires,
Un \[Am]c{\oe}ur, une Foi qui s'é\[G]lèvent,
Des \[C]Traces d'ave\[G]nir.
\endverse
\endsong

\beginsong{Allez viens}[by={Collège de branche, 2019}]
\beginchorus
\[Em]Allez \[C]viens, viens, \[G]viens dans la fo\[D]rêt !
\[Em]Allez \[C]viens, viens, \[G]viens dans la fo\[D]rêt !
\[Em]Bienvenue dans la peu\[C]plade : ton \[G]équipe, dans la ca\[D]bane.
\[Am]Mets ta che\[Em]mise, ton fou\[C]lard, al\[D]lez, re\[Em]joins-nous !
Allez \[Em]viens.
\endchorus
\beginverse
\[Em]Tu quittes pa\[Bm7]pa, ma\[C]man ; au re\[Bm]voir les pa\[Am7]rents.
\[Bm]Salut les co\[C]pains. En o\[D]range on est bien !
\[Em]Au c{\oe}ur de la na\[Bm7]ture, tu \[C]vas apprendre en \[Bm]jouant.
\[Am7]Tu verras, c'est pas si \[Bm]dur, tu \[C]viens vivre une aven\[D]ture !
\endverse
\beginverse
\[Em]Tu deviens louv'\[Bm7]teau jean\[C]nette ; jouons, dan\[Bm]sons, c'est la fête.
\[Am7]Après s'être pré\[Bm]sentés, on \[C]prend le goû\[D]ter.
\[Em]Tous ensemble pour \[Bm7]faire des n{\oe}\[C]uds, à plu\[Bm]sieurs c'est toujours mieux.
\[Am7]Des étoiles plein les \[Bm]yeux en \[C]chantant autour du \[D]feu.
\endverse
\beginverse
\[Em]Vas agir pour la na\[Bm7]ture et \[C]courir à toute al\[Bm]lure.
\[Am7]Tu vas grandir, dé\[Bm]couvrir la \[C]forêt, ses sen\[D]tiers.
\[Em]Dormir tous sous la \[Bm7]tente ou les \[C]étoiles, ça te \[Bm]tente.
\[Am7]Tu sauras tendre la \[Bm]main pour in\[C]viter des co\[D]pains.
\endverse
\beginverse
\[Em]Vas-y fais de ton \[Bm7]mieux, en\[C]semble c'est plus joy\[Bm]eux.
\[Am7]Maintenant c'est à ton \[Bm]tour, al\[C]lons voir aux alen\[D]tours.
\[Em]Les guider sur les che\[Bm7]mins, par\[C]tons en route vers de\[D]main.
\endverse
\endsong

\beginsong{Au cœur de la planète}[by={Hubert Bourel, Marie-Louise Valentin}]
\beginchorus
\[Dm]Un c{\oe}ur \[C]bat
\[A]Au c{\oe}ur de la pla\[Dm]nète, \[Dm]c'est ton \[C]c{\oe}ur,
\[A7]Louveteau Jean\[Dm]nette.
\[Gm]Tu t'habilles aux \[C7]couleurs du so\[Dm]leil le\[Gm7]vant,
\[Bm]Du so\[C]leil le\[Dm]vant.
\endchorus
\beginverse
\[Gm]La confiance est à l'ac\[Dm]cueil,
\[Gm]L'orée des forêts t'ap\[Dm]pelle.
\[Bm]L'amitié donne des \[F]ailes,
\[Gm6]Viens on va fran\[Gm7]chir le \[A]seuil...
\[Gm]Auprès de la peu\[Dm]plade,
\[Bm]Ce sont tes a\[F]mis,
\[Gm7]Tu vas réus\[A]sir ta \[Dm]vie.
\endverse
\beginverse
\[Gm]Pour l'ambiance enso\[Dm]leillée,
\[Gm]T'as des blagues plein les \[Dm]poches.
\[Bm]Des idées plein la ca\[F]boche,
\[Gm6]Des histoires \[Gm7]à racon\[A]ter...
\[Gm]Auprès de la peu\[Dm]plade,
\[Bm]Ce sont tes a\[F]mis,
\[Gm7]Tu vas réus\[A]sir ta \[Dm]vie.
\endverse
\beginverse
\[Gm]Tu vas devenir veil\[Dm]leur,
\[Gm]Tu grandis et tu t'en\[Dm]gages.
\[Bm]Tu signes au bas de la \[F]page,
\[Gm6]Le monde de\[Gm7]vient meil\[A]leur...
\[Gm]Auprès de la peu\[Dm]plade,
\[Bm]Ce sont tes a\[F]mis,
\[Gm7]Tu vas réus\[A]sir ta \[Dm]vie.
\endverse
\endsong

\beginsong{Promesse louveteaux jeannettes}[by={Hubert Bourel}]
\beginchorus
\[E]Dans ce monde qui a\[B]van\[E]ce
\[A]Je fais de mon \[B7]mieux
Pour \[E]faire vivre l'espé\[A]ran\[B]ce
Et \[E]grandir le \[B7]feu !
\endchorus
\beginverse
\[E]Je me connais, \[F#m]tu te connais
\[A]Je sais miser sur mes a\[B]touts
\[E]Je le promets, \[F#m]tu le promets
\[A]Ensemble, on ira jusqu'au \[B]bout.
\[C#m]Quand l'un de nous tombe\[F#m]ra
\[A]On pourra comp\[B7]ter sur toi.
\endverse
\beginverse
\[E]Je me connais, \[F#m]tu te connais
\[A]Tu sais qu'on sera toujours \[B]là
\[E]Je le promets, \[F#m]tu le promets
\[A]Même si j'oubliais la \[B]loi.
\[C#m]Quand l'un de nous tombe\[F#m]ra
\[A]On pourra comp\[B7]ter sur toi.
\endverse
\beginverse
\[E]Je me connais, \[F#m]tu te connais
\[A]Avec le secours du Sei\[B]gneur
\[E]Je le promets, \[F#m]tu le promets
\[A]Prenons la vie du côté \[B]cœur.
\[C#m]Quand l'un de nous tombe\[F#m]ra
\[A]On pourra comp\[B7]ter sur toi.
\endverse
\beginverse
\[E]Je me connais, \[F#m]tu te connais
\[A]Je sais prendre confiance en \[B]moi
\[E]Je le promets, \[F#m]tu le promets
\[A]La vie ouvrira grand ses \[B]bras.
\[C#m]Quand l'un de nous tombe\[F#m]ra
\[A]On pourra comp\[B7]ter sur toi.
\endverse
\beginverse
\[E]Je me connais, \[F#m]tu te connais
\[A]Je me sais déjà bâtis\[B]seur
\[E]Je le promets, \[F#m]tu le promets
\[A]Avec vous d'un monde meil\[B]leur.
\[C#m]Quand l'un de nous tombe\[F#m]ra
\[A]On pourra comp\[B7]ter sur toi.
\endverse
\endsong

\beginsong{Oser en grand}[by={Scouts et Guides de France}]
\beginchorus
\[F]Avec ma com\[C]mu, \[G]oser bâtir en \[Am]grand,
\[Fadd9]J'explore et j'ap\[C]prends, \[G]mon équipe m'at\[Am]tend.
\[F]Je veux changer le \[C]monde en \[G]rêvant sous la \[Am]toile,
\[F]Et vivre des odys\[C]sées, \[G]guidé par les é\[Am]toiles,
\[F]Avec \[C]toi. \[G] \[G7]
\endchorus
\beginverse
\[C]J'enfile ma chemise \[Em]bleue, \[Dm]direction le bon\[Am]heur,
\[F]En avoir plein les \[C]yeux, un \[F]monde tout en cou\[G]leur.
\[C]Écouter la na\[Em]ture, \[Dm]partir à l'aven\[Am]ture,
\[F]Je me lève et j'a\[C]gis, les \[F]défis d'aujourd'\[G]hui.
\endverse
\beginverse
\[C]Fier de ma che\[Em]mise, je \[Dm]défends mes va\[Am]leurs,
\[F]Construire des pro\[C]jets qui \[F]rayonnent dans nos \[G]c{\oe}urs.
\[C]Nos rêves et nos ta\[Em]lents \[Dm]pour la communau\[Am]té,
\[F]S'engager dans le \[C]monde \[F]pour le faire chan\[G]ger.
\endverse
\beginverse
\[C]Même sans ma che\[Em]mise, que \[Dm]partout l'on se \[Am]dise :
\[F]Ma vie est scintil\[C]lante car ma \[F]promesse est vi\[G]vante !
\[C]En famille, au col\[Em]lège \[Dm]c'est un grand privi\[Am]lège
\[F]D'avoir tant à don\[C]ner, des \[F]rires à parta\[G]ger.
\endverse
\beginverse
\[C]Je change de che\[Em]mise, \[Dm]destination ail\[Am]leurs,
\[F]Décoller vers mes \[C]rêves, \[F]transmettre mon bon\[G]heur.
\[C]Mais avant de par\[Em]tir \[Dm]Ami, je dois te \[Am]dire :
\[F]Ose, vivre et explo\[C]rer, \[F]à ton tour de jou\[G]er !
\endverse
\endsong

\beginsong{Avec ma tribu}[by={Benoît Gschwind}]
\beginchorus
\[G]Un monde tout en \[D]bleu,
Des \[G]projets plein les \[D]yeux,
Des \[G]terres à explo\[D]rer avec ma \[C]tri\[G]bu !
\[G]Une aventure en \[D]vrai,
À \[C]portée d'équi\[D]page,
Pour \[C]grandir chaque \[G]jour avec ma \[Em]tri\[D]bu ! \[G]
\endchorus
\beginverse
\[Em]Bienvenue à toi
\[Am]Qui \[Em]vient d'arriver,
\[C]Pour trouver et décou\[D]vrir ce qui nous rend \[G]heureux ! \[D7]
\[Em]Avec la tribu
\[Am]Nous \[Em]serons témoins
Des \[C]pas que tu \[D]feras d'aven\[G7]ture en aven\[D7]ture !
\endverse
\beginverse
\[Em]Un jour tu voudras
\[Am]Vivre \[Em]notre loi,
\[C]Avancer et t'enga\[D]ger avec ton équi\[G]page ! \[D7]
\[Em]Devant la tribu
\[Am]Tu \[Em]diras « je crois »,
Pro\[C]messe d'avan\[D]cer d'aven\[G7]ture en aven\[D7]ture !
\endverse
\beginverse
\[Em]Toujours partager
\[Am]Des ta\[Em]lents nouveaux,
\[C]S'accepter et s'en\[D]traider pour vivre en équi\[G]page ! \[D7]
\[Em]Et dans la tribu
\[Am]Un rôle \[Em]à chacun,
En\[C]semble nous i\[D]rons d'aven\[G7]ture en aven\[D7]ture !
\endverse
\beginverse
\[Em]Dieu prend nos chemins,
\[Am]Il est \[Em]notre ami,
\[C]Pain rompu et parta\[D]gé, il est le Dieu vi\[G]vant ! \[D7]
\[Em]Avec la tribu
\[Am]Prier \[Em]et chanter,
Ren\[C]contrer Jésus-\[D]Christ, d'aven\[G7]ture en aven\[D7]ture !
\endverse
\beginverse
\[Em]Avec les copains
\[Am]Rêver \[Em]à demain,
\[C]Explorer et inven\[D]ter : un défi pour cha\[G]cun ! \[D7]
\[Em]Avec la tribu
\[Am]Nous i\[Em]rons plus loin,
\[C]Camper et dé\[D]camper d'aven\[G7]ture en aven\[D7]ture !
\endverse
\endsong

\beginsong{Vis tes rêves}[by={Nicolas Mellier, Amplitude}]
\beginverse
\[Em]Des questions compli\[C]quées
Des \[G]avis parta\[D]gés. Moi au milieu des gens
\[Em]La Terre est effray\[C]ée
Les \[G]Hommes ont gaspil\[D]lé. Moi au milieu des gens
\[Em]FIER ! De savoir faire des \[C]choix, qui déplaisent par\[G]fois
Je refuse de me \[D]taire
\[Em]FIER ! De donner mes i\[C]dées, je construis des pro\[G]jets
Un avenir pour la \[D]Terre
\endverse
\beginchorus
\[Em]VIS TES \[C]RÊVES ! \[G]VIS TES \[D]RÊVES !
\[Em]Oser entrer dans la lu\[C]mière
\[G]Se lever, agir et être \[D]fier
\[Em]Rêver, avant que la \[C]nuit
N'\[G]s'achève
Rassem\[D]blés on vit nos Rêves !
\endchorus
\beginverse
\[Em]On reste mal éveil\[C]lé
De\[G]vant notre té\[D]lé. On s'évade sur une chaise
\[Em]Et quand tout est chif\[C]fré
Les \[G]Dollars, les i\[D]dées. J'ai pas de place pour le rêve
\[Em]J'RÊVE ! En couleur, en \[C]HD, d'un monde plus colo\[G]ré
Et plein de déci\[D]bels
\[Em]J'RÊVE ! Avec tous mes a\[C]mis, d'un futur où la \[G]vie
Chaque jour est plus \[D]belle
\endverse
\beginverse
\[Em]Des ombres qui se \[C]pressent
Des \[G]regards qui se \[D]baissent. Moi au milieu des gens
\[Em]Des Hommes comme des ma\[C]chines
Qui \[G]se suivent sur une \[D]ligne. Moi au milieu des gens
\[Em]J'OSE ! Interpeller les \[C]grands et sourire aux pas\[G]sants
Aussi tendre la \[D]main
\[Em]J'OSE ! Apprivoiser mes \[C]doutes, et puis prendre la \[G]route
Guidé par mon ins\[D]tinct
\endverse
\endsong

\beginsong{La caravane}[by={Amplitude}]
\beginverse
\[Am]Dans ce voyage aux mille visages
\[Am]Rencontrer l'autre sur la route
\[Am]Nos différences on les partage
\[Am]Partons ensemble sans bagages
\endverse
\beginchorus
\[F]La caravane prend le dé\[G]part
\[Am]Vers une desti\[F]nation bien \[C]con\[G]nue
\[Am]D'étapes en \[F]étapes, expé\[C]riences vé\[G]cues
\[Am]Sans toi il n'y \[F]aura pas d'his\[C]toi\[G]re
En \[Dsus2]route pour un nouveau voyage
\[F]Pour tracer ensemble notre sil\[G]lage
\endchorus
\beginverse
\[Am]Dans ce voyage aux mille visages
\[Am]Allez plus loin chaque seconde
\[Am]Prendre le temps, écrire une page
\[Am]Tendre sa main vers le monde
\endverse
\beginverse
\[Am]Dans ce voyage aux mille visages
\[Am]Oser le doute, les questions
\[Am]Vivre sa foi à chaque virage
\[Am]Pour affirmer ses passions
\endverse
\beginverse
\[Am]Dans ce voyage aux mille visages
\[Am]Poser son sac pour l'étape
\[Am]Vivre la fête, sous les étoiles
\[Am]Et regonfler son moral
\endverse
\beginverse
\[Am]Dans ce voyage aux mille visages
\[Am]Tenir le cap pour de bon
\[Am]Pour réussir notre passage
\[Am]Unissons-nous dans l'action
\endverse
\endsong

\beginsong{Je prends le relais}[by={Jérémie Dazy}]
\beginverse
\[Am]Malgré un monde qui se fissure
\[F]Les expériences me rassurent
\[C]M'ouvrir au monde, passer les murs
\[G]Rencontrer, vivre l'aventure.
\[Am]Les années passées à grandir
\[F]M'ont amené prêt pour servir
\[C]Les yeux tournés vers l'avenir
\[G]Un regard vert plein de plaisir
\endverse
\beginchorus
\[Am]Je prends le relais
\[F]Pour vivre ma \[G]vie, écrire des sourires
\[Am]On prend le relais
\[F]Pour agir en\[G]semble, là où bon nous semble
\[Am]Je suis, \[G]tu es, \[F]nous sommes
\[G]Des citoyens du monde
\[Am]Nous sommes, \[G]vous êtes, \[F]ils sont
\[Am]Des compa\[G]gnons
\endchorus
\beginverse
\[Am]Notre histoire s'écrit à plusieurs
\[F]Pour partager toutes nos valeurs
\[C]Les rires, l'espoir et les projets
\[G]Font de nous tous, une unité
\[Am]On a négocié des virages
\[F]On a avancé sans naufrage
\[C]Un seul mot d'ordre, fraternité
\[G]Pour une terre à partager
\endverse
\beginverse
\[Am]Au coin de toutes mes expériences
\[F]Je pose mon sac je fais silence
\[C]Je relis toutes mes actions
\[G]À quoi elles servent, vers qui elles vont
\[Am]Mais tout près au creux de mon cœur
\[F]Une promesse brille de mille lueurs
\[C]Ainsi commence le chemin
\[G]Qui me mènera vers demain
\endverse
\beginverse
\[Am]À l'heure de prendre son envol
\[F]Les souvenirs reprennent la main
\[C]Pour écrire l'histoire un peu folle
\[G]Des hommes et des femmes de demain
\[Am]Et quand il faudra s'engager
\[F]Avant qu'nos chemins se séparent
\[C]Jamais on n'oubliera, jamais
\[G]Cette chemise verte et ce foulard
\endverse
\endsong

\beginsong{Changer le monde}[by={Mark Hekic, Erck Ness}]
\beginverse
\[Dm]Même quand le ciel vire au gris,
Der\[Am]rière il cache un peu de bleu
\[C]Sous son armure mon cœur sourit,
\[G]Fais un effort découvre le
\[Dm]Si on te dit que la vie n'est pas belle,
Ré\[Am]ponds leur que c'est là tout l'enjeu
\[C]De voir dans les nuages des merveilles quand il \[G]pleut
\[Dm]Si on te dit que le monde est cruel,
Ré\[Am]ponds leur qu'on te le révèle mieux
\[C]On a tous un petit bout de soleil dans nos \[G]yeux
\endverse
\beginchorus
\[Dm]Puisqu'on est deux, qu'on se ras\[Am]semble
\[C]Puisqu'on le peut, qu'on se res\[G]semble
\[Dm]Impossible on le laisse au \[Am]vent
\[C]On est prêt, on est prêt, on est \[G]prêt
\[Dm]Rien ne sera jamais comme a\[Am]vant
\[C]Demain est un nouveau pré\[G]sent
\[Dm]Qu'on redessine comme des en\[Am]fants
\[C]On est prêt, on est prêt, on est \[G]prêt
À chan\[Dm]ger le monde, \[Am]monde
\[C]Oh, oh, oh, oh, \[G]oh
À chan\[Dm]ger le monde, \[Am]monde
\[C]Oh, oh, oh, oh, \[G]oh
\endchorus
\beginverse
\[Dm]On a oublié en chemin que nos vies
Sont des châ\[Am]teaux de sables
\[C]Dont il faut savoir prendre soin,
Non ce ne sont pas que des \[G]fables
\[Dm]Le bonheur c'est comme un boomerang
Quand \[Am]on le partage il nous revient
\[C]Si on a pas tous la même langue ça ne fait \[G]rien
\[Dm]C'est un point de vue, une question d'angle
De \[Am]verre moitié vide ou moitié plein
\[C]On a toujours plus que c'qu'il nous semble dans nos \[G]mains
\endverse
\beginverse
\[Dm]Et si la vie nous joue des tours
C'est \[Am]que c'est la règle du jeu
\[C]Mais on se relèvera tou\[G]jours
On fera de notre mieux
\[Dm]Et si la nuit remplace le \[Am]jour
Si tous nos rêves prennent \[F]feu
Ce qui nous sauvera c'est l'a\[G]mour
\endverse
\endsong

\beginsong{Lève-toi}[by={Jamboree}]
\beginverse
\[Am]C'est pas parce t'as comme un goût a\[F]mer, non
Au bord de tes lèvres des re\[C]lents de misère
C'est pas parce que la vie t'a ran\[G]gé dans un coin avec ton chagrin
\[Am]Qu'tu dois baisser les bras,
Ac\[F]cepter de te taire
\endverse
\beginverse
\[Am]C'est pas parce que ton cœur a pris la \[F]mer, non
Qu'la passion s'est noyée dans un \[C]tout petit verre
Les blessures invi\[G]sibles
\[Am]Elles te suivront partout, même si tu pars \[F]loin
Alors reviens au port, t'as des \[C]choses à y faire
\[Am]C'est pas parce que les gens t'ont dit qu'tu valais \[F]rien
Que c'est ton des\[C]tin
Alors reviens au port, t'as des \[G]bouches à faire taire
\endverse
\beginchorus
\[Am]Allez lève, lève, lève, é\[F]lève toi
\[C]Comme on lève le \[G]poing
\[Am]N'attends pas qu'on te donne \[F]raison ou pas
\[Am]Allez lève, lève, lève, é\[F]lève toi
\[C]Allez déplie tes rêves,
Re\[G]mue ta carcasse et bouge tes dix doigts
\endchorus
\beginverse
\[Am]C'est vrai que la vie est remplie de mys\[F]tères
Et que bien trop souvent on re\[C]garde en arrière
Mais plutôt qu'de pleurer sur c'que t'aurais pu \[G]faire
Ça ne sert à rien
\[Am]Pose la pierre, la pre\[F]mière
Du bonheur de de\[C]main
\endverse
\beginverse
\[Am]J'te promets pas qu'ton cœur n'aura plus ja\[F]mais mal, non
On n'se sépare jamais des cas\[C]seroles qu'on s'trimballe
Mais tu dois es\[G]sayer
\[Am]T'as vraiment rien à perdre, le bon\[F]heur c'est tout droit
Pour le saisir tu dois \[C]faire le premier \[G]pas
\endverse
\endsong

\beginsong{Comme un accord}[by={Davide Esposito, Laurent Lamarca}]
\textnote{Capo 2}
\beginverse
\[Em]On a tous un bout de \[C]nous à répa\[G]rer,
\[Em]Seul, on se voit flou et \[C]désempa\[G]ré.
\[Em]Sans les yeux des autres on se \[C]cherche en \[G]vain,
\[Em]Dans les yeux des autres il y a \[C]des che\[G]mins.
\endverse
\beginverse*
\[Em]Je n'suis pas la meilleure,
\[C]Je n'suis pas la plus \[G]forte.
\[Em]On n'est pas seul à bord,
\[C]On est beau à plu\[D]sieurs.
\endverse*
\beginchorus
Comme un ac\[G]cord,
Malgré la route et la rancœur,
\[Em]Le poids des doutes et les erreurs,
\[C]Les omissions et les remords,
\[Am]Être \[D]plus fort, comme un ac\[G]cord.
Malgré les soirs où on se perd,
\[Em]Les idées noires et la colère,
\[C]La peur du vide et les douleurs,
\[Am]Être \[D]meilleurs,
Comme un ac\[Em]cord.
\endchorus
\beginverse
\[Em]Qu'est-ce que je s'rais sans l'air \[C]que l'on m'a don\[G]né,
\[Em]Je veux souffler les mots, les \[C]renvo\[G]yer.
\[Em]J'entends dire que l'homme est un \[C]loup pour l'\[G]homme,
\[Em]Je veux croire en nous, en \[C]ce que nous \[G]sommes.
\endverse
\beginverse*
\[C]On a tous un \[D]bout de nous à répa\[Em]rer,
\[C]Seul, on se voit \[D]flou et désempa\[Bm]ré.
\endverse*
\endsong

\beginsong{Jamboree}[by={Ben Mazué, Davide Esposito, Renaud Rebillaud}]
\textnote{Capo 3}
\beginverse
\[C]T'as besoin d'aide, ça va pas ?
\[G]Tu veux que je reste avec toi ?
\[Am]Tu sais que tu pourras compter sur \[F]moi.
\[C]Je me souviens de nos promesses,
\[G]C'était pas rien ce qui nous reste,
\[Am]Tu sais que tu pourras compter sur \[F]moi.
\endverse
\beginverse*
\[C]Tout seul on ne reste qu'une goutte d'eau,
\[G]Soudé on devient le courant du ruisseau.
\[Am]Oui, c'est en étant nous-mêmes qu'on devient plus \[F]beau.
\endverse*
\beginchorus
\[C]Tu sais, ça paraît comme une évi\[G]dence,
Nous on se fout des appa\[Am]rences,
Si tu veux rentrer dans la \[F]danse.
\[C]Oh oui, tu sais, ça paraît comme une évi\[G]dence,
La force de nos diffé\[Am]rences,
Si tu veux rentrer dans la \[F]danse.
\endchorus
\beginchorus
\[C]Ohohoh, Jamboree, un peu d'amour et
\[G]Jamboree, on arrive, attend-nous,
\[Am]Jamboree, chanter avec vous,
\[F]La lala lala.
\endchorus
\beginverse
\[C]On ira marcher, mon ami,
\[G]Je compte sur toi, tu comptes sur lui.
\[Am]Partageons les sourires et les sou\[F]cis.
\[C]On se rapprochera tellement,
\[G]Moi le métal et toi l'aimant,
\[Am]Et si tu tombes, je tombe évidem\[F]ment.
\endverse
\endsong

\beginsong{Des couleurs sur mon chemin}[by={Nicolas Mellier, Amplitude}]
\beginverse
\[F#m]Tant de moments parta\[A]gés
Les \[D]beaux jours passés en\[A]sem\[B]ble
Des \[F#m]sourires par mil\[A]liers
Des \[D]délires qui nous res\[A]sem\[E]blent
\endverse
\beginverse
\[F#m]Des couleurs sur mon chemin
\[A]Des couleurs sur mon chemin
\[B]Des couleurs sur mon chemin
\[F#m]Des couleurs sur mon che\[A]min
\[F#m]Des couleurs sur mon chemin
\[A]Des couleurs sur mon chemin
\[B]Des couleurs sur mon chemin
\[F#m]Des couleurs sur mon che\[E]min
\endverse
\beginchorus
\[F#m]Une chemise un foulard \[D]et mon sac \[A]sur le \[C#m]dos
\[D]Ma pro\[A]messe, \[Bm]d'une main qui t'aide à passer la rivière
\[F#m]Une chemise un foulard \[D]et mon sac \[A]sur ton \[C#m]dos
\[D]Le scou\[A]tisme dans la \[C#]peau
\endchorus
\beginverse
\[F#m]On échange nos tré\[A]sors
Des \[D]rencontres sur la \[A]rou\[B]te
\[F#m]Une gourde au bout de l'ef\[A]fort
Les \[D]étoiles en clef de \[A]voû\[E]te
\endverse
\beginverse
\[F#m]Lorsque les épaules \[A]souffrent
\[D]On s'assoit près du che\[A]min\[B]
\[F#m]L'amour donne un deuxième \[A]souffle
\[D]Quand on partage le \[A]pain\[E]
\endverse
\beginverse
\[F#m]Les yeux plongés dans le \[A]ciel
\[D]Une tente et de l'air \[A]pur\[B]
\[F#m]Une cour sans fron\[A]tières
\[D]Généreuse est la na\[A]tu\[E]re
\endverse
\endsong

\beginsong{Le monsieur en chemise}[by={Marc Durand}]
\beginverse
\[Am]J'suis farfa, j'ai 6 ans
\[Am]J'vais aux scouts avec Maman
\[C]J'aime jouer, j'aime camper,
\[F]Car tout ça me fait rêver
\[C]Y'a un M'sieur qui m'regarde du haut de ses 40 ans
\[G]Qui me dit mon p'tit gars, « tu fais déjà comme les grands »
\endverse
\beginchorus
\[F]Et il avait \emph{clap clap} \[C]raison \emph{clap clap}
\[G]Le monsieur, le monsieur en che\[C]mise
\[F]Il avait \emph{clap clap} \[C]raison \emph{clap clap}
\[G]De me faire un p'tit peu la le\[C]çon
\endchorus
\beginverse
\[Am]J'suis orange, j'ai 8 ans
\[Am]Et j'ai peur sans ma maman,
\[C]J'sais rien faire, j'pleure tout le temps
\[F]Et j'me d'mande ce qui m'attend,
\[C]Y a un m'sieur qui m'regarde du haut de ses 19 ans
\[G]Qui me dit « t'inquiète pas tu vas kiffer pour longtemps »
\endverse
\beginverse
\[Am]Je suis bleu, j'ai 12 ans
\[Am]Et j'ai plus peur sans maman
\[C]Entouré de copains, ils ne sont jamais bien loin
\[F]J'ai appris, j'ai grandi, j'ai décou\[C]vert la vraie vie
\[G]Et l'monsieur qui me dit, « tu vas voir c'est pas fini »
\endverse
\beginverse
\[Am]Je suis rouge, j'ai 15 ans
\[Am]Et j'kiffe quand y a plus maman
\[C]Entouré des mêmes gens qu'y avait quand j'avais 8 ans
\[F]Y a un m'sieur qui m'regarde du \[C]haut de ses 25 ans
\[G]Qui me dit « mon p'tit gars t'es en train de dev'nir grand »
\endverse
\beginverse
\[Am]Je suis vert, j'suis majeur
\[Am]Et j'pense même plus à ma mère
\[C]À bâtir des projets, à partir à l'étranger
\[F]Y a un m'sieur qui me dit « après \[C]ça c'est terminé »
\[G]Je lui dis « sûr'ment pas à mon tour de faire rêver »
\endverse
\beginverse
\[Am]Je suis chef, j'suis acteur
\[Am]Je fais vivre plein de cœurs
\[C]Peu importe la couleur ou l'bafa d'animateur
\[F]Y a un ptit qui me dit est-ce que \[C]je peux m'amuser
\[G]Et j'lui dis sûr'ment pas va vendre des calendriers
\endverse
\beginchorus
\[F]Et maint'nant \emph{clap clap} \[C]c'est moi \emph{clap clap}
\[G]Le monsieur, le monsieur en che\[C]mise
\[F]Et maint'nant \emph{clap clap} \[C]c'est moi \emph{clap clap}
\[G]Qui m'permet de faire un p'tit peu la le\[C]çon
\endchorus
\beginverse
\[Am]J'suis violet, j'ai plus d'temps,
\[Am]Je suis scout je suis parent
\[C]J'organise des soirées et des réu tout le temps
\[F]Y a un m'sieur qui me dit est ce que \[C]je peux m'engager
\[G]Et j'lui dis pourquoi pas tu peux devenir RG
\endverse
\beginchorus
\[F]Et maintenant \emph{clap clap} c'est \[C]lui \emph{clap clap}
\[G]Le monsieur le monsieur en che\[C]mise
\[F]Et maintenant \emph{clap clap} c'est \[C]lui \emph{clap clap}
\[G]Qui s'permet de faire un peu la le\[C]çon
\endchorus
\endsong

\beginsong{Chant de la promesse traditionnel}[by={Père Jacques Sevin}]
\beginverse
\[D]Devant tous je m'engage
\[Bm]Sur \[E7]mon hon\[A]neur,
Et \[D]je te fais hommage
\[G]De \[A]moi, Sei\[D]gneur !
\endverse
\beginchorus
\[G]Je veux t'aimer sans \[D]cesse,
\[Em]De \[A]plus en \[D]plus,
\[G]Protège ma pro\[D]messe,
\[G]Sei\[A]gneur Jé\[D]sus !
\endchorus
\beginverse
\[D]Je jure de te suivre
\[Bm]En \[E7]fier chré\[A]tien,
Et \[D]tout entier je livre
\[G]Mon \[A]cœur au \[D]tien !
\endverse
\beginverse
\[D]Fidèle à ma Patrie
\[Bm]Je \[E7]le se\[A]rai ;
\[D]Tous les jours de ma vie,
\[G]Je \[A]servi\[D]rai.
\endverse
\beginverse
\[D]Je suis de tes apôtres
\[Bm]Et \[E7]chaque \[A]jour
\[D]Je veux aider les autres
\[G]Pour \[A]ton a\[D]mour.
\endverse
\beginverse
\[D]Ta règle a sur nous-mêmes
\[Bm]Un \[E7]droit sa\[A]cré :
\[D]Je suis faible, tu m'aimes,
\[G]Je \[A]maintien\[D]drai.
\endverse
\endsong

\beginsong{Chant de la promesse}[by={Emmanuelle Audras, Philippe Bancon}]
\beginverse
\[C]Aujourd'hui je me sens \[F]prêt
\[C]Près de toi pour t'écou\[F]ter
\[C]Mon pas n'est pas assu\[F]ré
\[C]Mais nous sommes à tes cô\[G]tés
\[C]Ma parole est enga\[F]gée
\[C]Aux nôtres elle s'est mélan\[F]gée
\[C]Aujourd'hui je me sens \[G]prêt
\endverse
\beginchorus
\[F]Être utile avec mes \[C]sœurs
\[F]Donner de moi le meil\[C]leur
\[F]Oh Seigneur par ma pro\[C]messe
\[G]Je veux a\[C]gir
\endchorus
\beginverse
\[C]Je vous confie ma pa\[F]role
\[C]Nous te confions notre \[F]loi
\[C]Je la prends comme bous\[F]sole
\[C]Elle saura guider nos \[G]pas
\[C]Si nos pieds trébuchent au \[F]sol
\[C]Nous pourrons compter sur \[F]toi
\[C]Je vous confie ma pa\[G]role
\endverse
\beginchorus
\[F]Partager avec mes \[C]frères
\[F]Quels que soient les vents con\[C]traires
\[F]Oh Seigneur par ma pro\[C]messe
\[G]Je veux ser\[C]vir
\endchorus
\beginverse
\[C]Je sens la vie qui m'ap\[F]pelle
\[C]Elle sourit aux auda\[F]cieux
\[C]Je veux déployer mes \[F]ailes
\[C]Sers le monde il est pré\[G]cieux
\[C]Je veux être une étin\[F]celle
\[C]Le bonheur est conta\[F]gieux
\[C]Je sens la vie qui m'ap\[G]pelle
\endverse
\beginchorus
\[F]Prendre soin des plus fra\[C]giles
\[F]Je suis de la même ar\[C]gile
\[F]Oh Seigneur dans ma pro\[C]messe
\[G]Je te re\[C]çois
\endchorus
\endsong

\beginsong{Ensemble on est mieux}[by={Théophile Rénier}]
\textnote{Capo 1}
\beginverse
\[Em]J'm'en souviens, j'avais 6 ans
\[Am]J'comprenais pas c'que voulaient mes parents
\[C]Quand ils m'ont dit « Mon enfant,
\[B7]Va chez les scouts, c'est important ! »
\[Em]Moi, j'flippais grave au début
\[Am]J'sortais d'chez moi, d'mes habitudes
\[C]Mais la ribambelle m'ouvre la porte
\[B7]Et au final le rire m'emporte
\endverse
\beginchorus
\[Em]Ensemble on est mieux
\[C]On a du mal à s'dire adieu
\[G]Les scouts nous portent, nous transportent
\[B7]Nous font danser comme le feu
\[Em]Ensemble on est mieux
\[C]On a du mal à s'dire adieu
\[G]Les scouts nous portent, nous transportent
\[B7]Nous font danser comme le feu
\endchorus
\beginverse
\[Em]Puis j'me rappelle de mes 8 ans
\[Am]Chez les Louv'teaux tout était différent
\[C]Je découvrais la vie au grand air
\[B7]Les jeux dans l'bois, la vie sans manières
\[Em]On court partout, on fait les fous
\[Am]On rit, on chante, comme des loups
\[C]Avec la meute j'deviens plus grande
La force du loup, c'est le clan
\endverse
\beginverse
\[Em]Arrivé chez les Éclaireurs
\[Am]Plus personne ne me faisait peur
\[C]Quand j'suis en bande avec mes potes
\[B7]J'me sens plus fort même pour la tot'
\[Em]On vit ensemble et en patrouille
\[Am]On devient les rois d'la débrouille
\[C]On construit même nos pilotis
\[B7]À bout de bras et d'énergie
\endverse
\beginverse
\[Em]Puis viennent les pi's et leurs envies
\[Am]D'aventure et de fantaisie
\[C]Pour ça on rêve toute l'année
\[B7]On veut qu'une chose c'est tout changer
\[Em]Durant le camp on s'ouvre à tout
\[Am]Une aut' culture, un aut' chez nous
\[C]On vient, on aide, si on peut
\[B7]On échange même avec des vieux
\endverse
\beginverse
\[Em]Enfin, tu es animateur
\[Am]Du temps, du talent et du cœur
\[C]Et tu t'engages bénévol'ment
\[B7]À transmettre tout ce que t'as dans l'sang
\[Em]Parfois c'est vrai, c'est la galère
\[Am]Mais c'est pas grave, en staff tu gères
\[C]Baden-Powell est fier de toi
\[B7]Chante avec moi et lève 3 doigts
\endverse
\endsong

\beginsong{Viens mélanger tes couleurs}[by={Moulin Pélissier}]
\beginchorus
\[D]Viens mélanger tes cou\[A]leurs avec \[Bm]moi
Réveiller le bon\[G]heur qui \[A]dort au fond de \[D]toi
Faire jaillir la lu\[A]mière de nos \[Bm]vies
Improviser la \[G]fête au plein \[A]cœur de la \[D]nuit
\endchorus
\beginverse
\[G]Tu connais la mi\[A]sère
\[D]Qui condamne au si\[Bm]lence
\[G]Prends la main de tes \[A]frères
\[D]Invente un pas de \[G]dan\[A]se.
\endverse
\beginverse
\[G]Tu réclames en par\[A]tage
\[D]Ce qui permet la \[Bm]vie
\[G]Relève ton vi\[A]sage
\[D]Vois l'étoile dans la \[G]nuit.\[A]
\endverse
\beginverse
\[G]Tu t'opposes à la \[A]force
\[D]Qui tue la liber\[Bm]té
\[G]Entrouvre ton é\[A]corce
\[D]Au soleil de l'é\[G]té.\[A]
\endverse
\beginverse
\[G]Tu brandis la ban\[A]nière
\[D]De la fraterni\[Bm]té
\[G]Dépasse les fron\[A]tières
\[D]En apportant la \[G]paix.\[A]
\endverse
\endsong

\beginsong{La promesse}[by={Grégoire, Jean-Jacques Goldman}]
\textnote{Capo 3}
\beginverse
\[Em]On était quelques \[C]hommes
Quelques \[Am]hommes, quelques femmes \[D]rêvant de liberté
\[Em]On n'était pas à \[C]vendre
Mais on \[Am]pouvait revendre des \[D]montagnes d'ami\[G]tié
Le cœur en bandou\[Em]lière
\[C]Et les bras grands ouverts à \[Am]tous les étran\[D]gers
On n'avait pas de \[Em]peur
\[C]On sentait la chaleur qu'on \[Am]savait se don\[B7]ner
\endverse
\beginverse
\[Am]Même au fin fond du dé\[Em]sert
On \[D]aidait les plus faibles à ne \[G]jamais tom\[Am]ber
Même au milieu des chi\[Em]mères
\[C]On y croyait plus fort quand le \[B7]courage manquait
\endverse
\beginchorus
\[G]Oh vous mes compagnons mes amis de jeu\[D]nesse
\[Em]Quelles que soient vos histoires non n'ou\[B7]bliez jamais
\[G]Qu'un beau jour nous avions fait ensemble une pro\[D]messe
\[C]S'il n'en reste qu'un nous serons ce der\[B7]nier
\endchorus
\beginverse
\[Em]On était plein d'ar\[C]deur
Et on \[Am]sortait vainqueur de nos \[D]pauvres bles\[Em]sures
Quand les pleurs étaient \[C]lourds
On se \[Am]trouvait toujours une \[D]voix qui nous ras\[G]sure
On avait tant d'en\[Em]vie
\[C]Qu'on voyait notre vie comme une \[Am]belle aven\[D]ture
On n'avait pas de \[Em]maître
\[C]La seule à nous soumettre était la \[Am]mère na\[B7]ture
\endverse
\beginverse
\[Am]Même au fin fond du dé\[Em]sert
On \[D]aidait les plus faibles et \[G]quitte à y res\[Am]ter
Même au milieu des chi\[Em]mères
\[C]On y croyait plus fort quand le \[B7]courage manquait
\endverse
\beginchorus
\[G]Oh vous mes compagnons mes amis de jeu\[D]nesse
\[Em]Quelque soit mon histoire non je n'ou\[B7]blie jamais
\[G]Et aujourd'hui encore je refais la pro\[D]messe
\[C]S'il n'en reste qu'un nous serons ce der\[B7]nier
\[G]Oh vous mes compagnons mes amis de jeu\[D]nesse
\[Em]Quelles que soient vos histoires, ne m'ou\[B7]bliez jamais
\[G]Et si un jour je tombe faites moi cette pro\[D]messe
\[C]S'il n'en reste qu'un vous serez ce der\[B7]nier
\[Am]Ce der\[B7]nier, ce der\[Em]nier
\endchorus
\endsong

\beginsong{Promesse}[by={E. Audras, P. Bancon}]
\beginverse
\[C]Aujourd'hui je me sens \[F]prêt.
\[C]Mon pas n'est pas assu\[G]ré.
\[C]Ma parole est enga\[F]gée.
\[C]Aujourd'hui je me sens \[G]prêt.
\endverse
\beginchorus
\[C]Être utile avec mes \[F]s{\oe}urs,
\[C]Donner de moi le meil\[F]leur.
\[C]Oh Seigneur par ma pro\[F]messe,
\[G]Je veux a\[C]gir !
\endchorus
\beginverse
\[C]Je vous confie ma pa\[F]role.
\[C]Je la prends comme bous\[G]sole.
\[C]Si nos pieds trébuchent au \[F]sol,
\[C]Je vous confie ma pa\[G]role.
\endverse
\beginverse
\[C]Je sens la vie qui m'ap\[F]pelle.
\[C]Je veux déployer mes \[G]ailes.
\[C]Je veux être une étin\[F]celle.
\[C]Je sens la vie qui m'ap\[G]pelle.
\endverse
\endsong

\beginsong{Red River Valley}
\beginverse
\[G]Les pionniers sont passés \[C]avant le \[G]jour
Dans les rues du village acca\[D7]blé
\[G]Et mon cœur a \[G7]frémi à leur \[C]pas lourd
Sur les \[G]bords de la \[D]Red River \[G]Valley
\endverse
\beginchorus
\[G]O Seigneur la roue \[C]tourne entre tes \[G]mains
Où je vais aujourd'hui je ne \[D]sais
\[G]O Seigneur la roue \[G7]tourne entre tes \[C]mains
Mais je \[G]veux retrou\[D]ver les pion\[G]niers
\endchorus
\beginverse
\[G]Les pionniers ont pei\[C]né pour le \[G]village
À leurs mains la vallée s'est pli\[D7]ée
\[G]Et mes yeux ont \[G7]vu naître un \[C]barrage
Sur les \[G]bords de la \[D]Red River \[G]Valley
\endverse
\beginverse
\[G]Les pionniers ont mar\[C]qué dans la \[G]clairière
Que le pain se partage entre \[D7]tous
\[G]Et ma main s'est ou\[G7]verte à mes \[C]frères
Sur les \[G]bords de la \[D]Red River \[G]Valley
\endverse
\beginverse
\[G]Les pionniers ont chan\[C]té dans la \[G]nuit claire
Que la terre est à qui la vou\[D7]lait
\[G]Et ma voix s'est u\[G7]nie à leur \[C]chant fier,
Sur les \[G]bords de la \[D]Red River \[G]Valley
\endverse
\beginverse
\[G]Les pionniers ont pro\[C]mis de re\[G]venir
L'herbe pousse aujourd'hui à nos \[D7]pieds
\[G]Et mon cœur s'est trou\[G7]vé fait pour \[C]servir
Sur les \[G]bords de la \[D]Red River \[G]Valley
\endverse
\endsong

\beginsong{Des signes de fraternité}[by={Mannick, Laurent Grzybowski}]
\textnote{Capo 2}
\beginverse
\[Em]Ça chauffe un peu partout sur la pla\[C]nète, \[D] \[Em]
\[Em]Au c{\oe}ur des Mégapoles, on se \[C]flingue à tout \[D]va. \[Em]
\[Em]La Raison du plus fort est la plus \[C]bête, \[D] \[Em]
\[Em]Mais par-dessus tout ça, \[C]regarde bien tu ver\[D]ras :
\endverse
\beginchorus
\[C]Des signes de Frater\[D]ni\[G]té,
\[Em]Dans tous les re\[C]coins de la \[G]terre. \[D]
\[C]Des signes de frater\[D]ni\[Em]té,
\[Am]Sur les barbelés de la \[B7]guerre.
\[C]Des Signes de Frater\[D]ni\[G]té.
\[Em]Des \[C]signes de Fra\[D]terni\[Em]té.
\endchorus
\beginverse
\[Em]Toujours chacun pour soi, chacun sa \[C]route, \[D] \[Em]
\[Em]On fait des hécatombes au volant de la \[C]vie. \[D] \[Em]
\[Em]Pour passer le premier quoi qu'il en \[C]coûte, \[D] \[Em]
\[Em]Mais par-dessus tout ça, \[C]regarde bien tu ver\[D]ras :
\endverse
\beginverse
\[Em]Partout les vieux démons reviennent en \[C]force, \[D] \[Em]
\[Em]Et braillent des slogans pour bannir l'étran\[C]ger. \[D] \[Em]
\[Em]D'un quartier d'un pays, au c{\oe}ur des \[C]hommes, \[D] \[Em]
\[Em]Mais par-dessus tout ça, \[C]regarde bien tu ver\[D]ras :
\endverse
\beginverse
\[Em]Des peuples cousus dans la mi\[C]sère \[D] \[Em]
\[Em]Et croisent l'abondance aux écrans des nan\[C]tis. \[D] \[Em]
\[Em]De quoi désespérer la vie en\[C]tière, \[D] \[Em]
\[Em]Mais par-dessus tout ça, \[C]regarde bien tu ver\[D]ras :
\endverse
\beginverse
\[Em]Devant la religion du terro\[C]risme, \[D] \[Em]
\[Em]Barbus de Barbarie, fous d'un dieu-Atti\[C]la. \[D] \[Em]
\[Em]Le monde a des relents d'appoca\[C]lypse, \[D] \[Em]
\[Em]Mais par-dessus tout ça, \[C]regarde bien tu ver\[D]ras :
\endverse
\beginverse
\[Em]Regarde tu verras, le monde \[C]change, \[D] \[Em]
\[Em]Il y a des mains tendues et des c{\oe}urs grands ou\[C]verts \[D] \[Em]
\[Em]Qui font se mélanger les diffé\[C]rences, \[D] \[Em]
\[Em]Dans les rues de ta vie, tu vois fleurir aujourd'\[C]hui ! \[D] \[Em]
\endverse
\endsong

\beginsong{Comme un appel}[by={E. Audras, P. Bancon / Amplitude}]
\beginchorus
\[G]Comme un ap\[D]pel résonne en \[Em]toi
Parmi le \[C]flot des « comment » des « pour\[G]quoi »,
Sauter le \[D]pas et entrer dans la \[Em]danse. \[F]
Cet ap\[C]pel comme une \[G]chance. \[Em]
\endchorus
\beginverse
\[Em]Du nord au sud, des\[D]sus dessous on me \[C]pille,
Je \[Em]suis sur deux bé\[D]quilles. \[C] \[Em]
\[Em]Ton oxy\[D]gène, ton \[C]sel, ton sol, ton \[Em]eau claire,
Je \[D]suis ta \[C]terre ! \[Bm] \[B]
\endverse
\beginverse
\[Em]Éclat de rêve dans le sou\[D]rire d'un en\[C]fant,
\[Em]Grandir chemin fai\[D]sant. \[C] \[Em]
\[Em]Cris d'injustice qui re\[D]fusent de se \[C]taire,
Je \[D]suis ton \[C]frère ! \[Bm] \[B]
\endverse
\beginverse
\[Em]Caché au creux du vent, d'un \[D]souffle lé\[C]ger,
\[Em]Je suis ta liber\[D]té. \[C] \[Em]
\[Em]Dans la tempête ma pa\[D]role est un \[C]feu,
Je \[D]suis ton \[C]Dieu ! \[Bm] \[B]
\endverse
\endsong

\beginsong{Le monde m'appelle}[by={Marydo Pellarini, Philippe Bancon}]
\textnote{Capo 2}
\beginverse
\[G]Il faut que j'aille voir ail\[C]leurs
Vers d'autres \[G]cieux et d'autres \[D]fleurs
Et \[G]s'il n'y a pas de che\[C]min
Je saurai \[G]bien tracer le \[D]mien
\[Em]Je pars, oui, mais je n'ou\[C]blie rien
De mon \[D]passé ni de mes \[G]liens
Et \[Em]j'ai glissé dans mes ba\[C]gages
Bien \[D]plus que des i\[G]mages
\endverse
\beginchorus
Le \[C]monde m'appelle !
\[C]Le monde m'at\[D]tend !
\[D]La vie est si \[C]belle
Je ne \[D]veux pas perdre mon \[G]temps
\[D]Le \[C]monde m'appelle !
\[C]Le monde m'at\[D]tend !
\[D]La vie est si \[C]belle
Je ne \[D]veux pas perdre mon \[G]temps\[D]
\endchorus
\beginverse
\[G]Je ne peux pas me conten\[C]ter
De la \[G]maison ou du ly\[D]cée
\[G]Tous les autres ont tant à m'ap\[C]prendre
Ils m'aide\[G]ront à mieux m'com\[D]prendre
\[Em]Je ne veux pas baisser les \[C]bras
Ail\[D]leurs on a besoin de \[G]moi
\[Em]Même si parfois c'est diffi\[C]cile
Je \[D]sais que je peux être u\[G]tile
\endverse
\beginverse
\[G]Non, je ne veux pas reje\[C]ter
Ce que \[G]mes parents m'ont don\[D]né
\[G]Mais s'ils m'ont offert la \[C]vie
C'est pour \[G]la vivre comme j'ai en\[D]vie
\[Em]Et je construirai ma mai\[C]son
Diffé\[D]rente de leur mai\[G]son
\[Em]Je leur emprunterai des \[C]pierres
Mais la \[D]ferai à ma ma\[G]nière
\endverse
\endsong

\beginsong{Promesse de scout}[by={Le Cinquième}]
\textnote{Capo 1}
\beginverse
\[Dm]Chez les Scouts, j'y passe le plus clair de mon temps,
\[G]En réunion entre potes, en week-end ou au camp.
\[Am]Nous on kiffe, on s'en fout du regard des gens,
\[F]Moi je suis scout un jour et heureux tout le temps !
\[Dm]Sapé comme jamais, ma chemise me rend cool,
\[G]Même s'il manque un insigne et qu'elle est toujours en boule.
\[Am]Le foulard autour du cou, toujours bien noué,
\[F]Car j'ai maté un tuto pour faire le n{\oe}ud carré.
\endverse
\beginverse
\[Dm]Avec mon équipe, monter la tente me rend fier,
\[G]Même avec un bout de bois pour caler la faîtière.
\[Am]Je construis des tables, je construis des bancs,
\[F]Mais une fois sur deux ils ne tiennent pas jusqu'à la fin du camp.
\[Dm]Désormais plus besoin d'aller à la salle,
\[G]Je me muscle en portant les jerricans et les malles.
\[Am]Passer une journée à creuser des feuillées,
\[F]Même si au final tout le monde va pisser dans la forêt.
\endverse
\beginverse
\[Dm]J'ai fait des gamelles, des thèques, des sioules,
\[G]J'ai chanté à tue-tête la scoutance et le matou.
\[Am]Quand je repense à tous les souvenirs que j'ai avec vous,
\[F]Je ne pense pas à moi, mais je pense à nous.
\endverse
\beginverse*
\[Dm]Je parle de souvenirs,
\[G]Mais si t'es pas scout tu peux pas connaître.
\[Am]Je parle de souvenirs,
\[F]Mais si t'es pas scout...
\endverse*
\beginchorus
\[Dm]J'vais repartir en camp, \[G]revenir tous les ans,
\[Am]Même quand je serai grand, \[F]promesse de Scout.
\endchorus
\beginverse
\[Dm]Parfois j'ai pas de chance pour la météo,
\[G]Concu sous la flotte, canicule en explo.
\[Am]Je garde le sourire, même s'il ne fait pas chaud,
\[F]J'allume un p'tit feu, j'sors les chamallows.
\[Dm]Je voudrais dormir, je suis épuisé,
\[G]Car depuis 6h du mat' les louveteaux sont réveillés.
\[Am]Pour les imaginaires, je sais improviser,
\[F]Un drap et un sac poubelle et me voilà déguisé.
\endverse
\beginverse
\[Dm]J'connais le chant de la Promesse par c{\oe}ur,
\[G]Mes calendriers trouvent toujours preneurs.
\[Am]J'ai inventé un jeu pour une réunion,
\[F]C'est un poule renard vipère où j'ai juste changé les noms.
\[Dm]Ici on mange mieux qu'au restau,
\[G]On mélange du éco+ avec des produits locaux.
\[Am]On cuisine tout au feu de bois,
\[F]Petit bonus : j'ai d'la cendre dans tous mes plats.
\endverse
\beginverse
\[Dm]Essayer de sortir sans se faire griller la nuit,
\[G]Sentir le feu de bois pendant la moitié de ma vie.
\[Am]Pas réussir à mettre tout dans mon sac à la fin du camp,
\[F]Être nostalgique de ces souvenirs, de ces moments.
\endverse
\beginverse*
\[Dm]J'pensais que pendant le camp j'avais bronzé,
\[G](Le monde m'appelle)
\[Am]Mais après 2 douches la crasse s'est enlevée,
\[F](Le monde m'attend)
\endverse*
\beginverse*
\[Dm]Une chemise, un foulard et mon sac sur le dos,
\[G]Ma Promesse, d'une main qui t'aide à passer la rivière.
\[Am]J'ai appris, j'ai grandi, j'ai découvert la vraie vie,
\[F]Ensemble on va plus loin, tu connais déjà le refrain.
\endverse*
\endsong

\beginsong{La scoutance}[by={Le Cinquième}]
\beginverse
\[D#m]J'vais te faire kiffer tes week-ends mon pe\[F#]tit
Me parle plus de tes profs, oublie tes poé\[C#]sies
Toutes les 3 semaines, ce sera l'eupho\[B]rie
Rendez-vous au local avec tous tes amis
\[D#m]Depuis le plus jeune âge, on a fait les 400 \[F#]coups
Bientôt dans l'calendrier, y'aura des photos de \[C#]nous
Thomson le vieux fermier, a viré son ma\[B]tou
Toujours éliminer les mêmes au 1er tour du loup-garou
\endverse
\beginverse
\[D#m]Voyager jusqu'au Nirva\[F#]na
Hôtel 1000 étoiles, dans mon duvet il fait pas \[C#]froid
Noter mon nom sur mon matos Que\[B]chua
Gober des flambys, manger des bananes-chocolat
\endverse
\beginchorus
\[D#m]Plus d'différences, allez rentre dans la \[F#]danse
C'est la scoutance, c'est la scou\[C#]tance
D'la bienveillance, délivrée sans ordon\[B]nance
C'est la scoutance, c'est la scoutance
\endchorus
\beginverse
\[D#m]Pas d'Uber Eats, ici on cuisine sur le \[F#]feu
Même quand il n'y a pas de bois, quand il fait froid ou quand il \[C#]pleut
On a tous promis qu'on ferait de notre \[B]mieux
Mais pour plier les tentes, je ne suis pas si consciencieux
\[D#m]Pour les installations, on se sert de scies, des ha\[F#]chettes
Et si ça tourne mal, on sait qu'on a 10\% de \[C#]pertes
Au camp les plus grands m'apprennent plein de techniques se\[B]crètes
Comme comment éviter la vaisselle de la tartiflette
\endverse
\beginverse
\[D#m]Quand je lis à la messe, je fais des étin\[F#]celles
Et avec j'ai rallumé la lumière de Bethle\[C#]em
Pendant toutes les veillées, c'est vrai, je me dé\[B]chaîne
Mais le sommeil ne viendra qu'après la fin du 5ème
\[D#m]J'vais faire un tour, j'ai mis mes chaussures de ran\[F#]do
Ma chemise mon foulard, et mon sac sur le \[C#]dos
Après toutes ces années, je suis devenu mon propre hé\[B]ros
Et quand je serai grand, je passerai le flambeau
\endverse
\endsong

\beginsong{Clameurs du monde}[by={François Tardif, Johan Tardif, Maïwenn Tardif}]
\beginverse
\[Bm]Malgré les peurs, les tem\[G]pêtes, les appels de la pla\[D]nète,
Un \[A]cri per\[Bm]siste.
Sous les cendres des \[G]guerres, les droits du monde ré\[D]sistent,
Nos \[A]c{\oe}urs gran\[Bm]dissent.
Que nos clameurs fran\[G]chissent les murs et les fron\[D]tières,
Que les \[A]voix des peuples s'u\[Bm]nissent.
Clameurs de ceux qu'on ou\[G]blie, de ceux qui es\[D]pèrent,
Et pour \[A]les démunis, rendons jus\[Bm]tice.
\endverse
\beginverse*
En\[G]semble, on ré\[D]pond à l'ap\[A]pel du \[Bm]monde,
En\[G]semble, on ré\[D]pond à l'ap\[A]pel.
\endverse*
\beginchorus
\[Bm]Écoute, les \[G]voix rugir du \[D]fond de nos \[A]c{\oe}urs !
\[Bm]Écoute, les \[G]voix rugir \[D]oh, oh ! \[A]
Faisons monter nos \[Bm]clameurs, \[G]clameurs, \[D]clameurs oh, \[A]oh !
Faisons monter nos \[Bm]clameurs, \[G]clameurs, \[D]clameurs oh, \[A]oh !
\endchorus
\beginverse
\[Bm]Malgré les boulever\[G]sements, les changements de la \[D]terre,
La \[A]joie appa\[Bm]raît.
Sous le regard silen\[G]cieux, toutes les nations s'af\[D]frontent,
La \[A]paix re\[Bm]naît.
Que nos clameurs dé\[G]voilent nos espoirs et nos \[D]rêves,
Que \[A]toutes nos angoisses dispa\[Bm]raissent.
Clameurs des c{\oe}urs bri\[G]sés, des âmes déchi\[D]rées,
Et pour \[A]l'avenir, une pro\[Bm]messe.
\endverse
\endsong

\beginsong{Des milliards de chemins}[by={Laurent Grzybowski}]
\beginchorus
\[Em]Des milliards de che\[C]mins,
\[D]Mais un seul pour cha\[Em]cun.
\[Em]Chemin d'amour qui fait gran\[C]dir,
\[D]Chemin qui trace un ave\[Em]nir.
\endchorus
\beginverse
\[G]Ton aventure est une \[D]marche
\[Am]Dans une troupe au c{\oe}ur bat\[Em]tant.
\[C]Auprès de toi je suis par\[G]tant,
\[Am]Comme un ami qui t'accom\[D]pagne.
\[C]Ensemble \[D]nous cher\[G]chons
\[C]Des routes \[G]qui li\[D]bèrent,
\[C]Des routes \[Am]de vi\[D]vants !
\endverse
\beginverse
\[G]Tu as des rêves plein la \[D]tête
\[Am]Et des projets qui portent \[Em]loin.
\[C]Auprès de toi je suis quel\[G]qu'un
\[Am]Qui veut comprendre ta jeu\[D]nesse.
\[C]Ensemble \[D]nous al\[G]lons
\[C]Vers l'aube \[G]et ses lu\[D]mières,
\[C]Vers l'aube \[Am]des vi\[D]vants !
\endverse
\beginverse
\[G]Tu connaîtras le froid des \[D]doutes,
\[Am]Les heures sombres de la \[Em]nuit.
\[C]Au plus profond j'entends ton \[G]cri
\[Am]Et je demeure à ton é\[D]coute.
\[C]Ensemble \[D]nous par\[G]lons
\[C]Des terres \[G]d'espé\[D]rances,
\[C]Des terres \[Am]de vi\[D]vants !
\endverse
\beginverse
\[G]Tu partiras au vent du \[D]large,
\[Am]Malgré la peur de l'oura\[Em]gan.
\[C]Je dis ma foi au Dieu pré\[G]sent,
\[Am]Le Christ appelle à des au\[D]daces.
\[C]Ensemble \[D]nous pre\[G]nons
\[C]La barque \[G]du pas\[D]sage,
\[C]La barque \[Am]des vi\[D]vants !
\endverse
\endsong

\beginsong{La scoutitude}
\beginverse
% TODO: ajouter paroles et accords
\endverse
\endsong

\beginsong{La strasbourgeoise}
\textnote{Capo 1}
\beginverse
\[Am]Petit papa, voici la mi-carême,
\[Dm]Car te voici déguisé en sol\[Am]dat.
\[Dm]Petit papa, dis-moi si c'est pour \[Am]rire,
Ou pour \[E]faire peur aux tous petits en\[Am]fants.
\endverse
\beginverse
\[Am]Non mon enfant, je pars pour la Patrie,
\[Dm]C'est un devoir où tous les papas s'en \[Am]vont.
\[Dm]Embrasse-moi, petite fille ché\[Am]rie,
Je ren\[E]trerais bien vite à la mai\[Am]son.
\endverse
\beginverse
\[Am]La neige tombe aux portes de la ville,
\[Dm]Là est assise une enfant de Stras\[Am]bourg.
\[Dm]Elle reste là malgré le froid, la \[Am]bise,
Elle reste \[E]là malgré le froid du \[Am]jour.
\endverse
\beginverse
\[Am]Gardez votre or, je garde ma puissance,
\[Dm]Soldat prussien, passez votre che\[Am]min.
\[Dm]Moi je ne suis qu'une enfant de la \[Am]France,
À l'enne\[E]mi je ne tends pas la \[Am]main.
\endverse
\beginverse
\[Am]Vous avez eu l'Alsace et la Lorraine,
\[Dm]Vous avez eu des millions d'étran\[Am]gers.
\[Dm]Vous avez eu Germanie et Bo\[Am]hème,
Mais mon p'\[E]tit c{\oe}ur il restera fran\[Am]çais.
\endverse
\endsong

\beginsong{Le retour à la terre}[by={Les Fatals Picards}]
\textnote{Capo 10}
\beginverse
\[Am]Minou fais tes valises et les miennes aus\[G]si,
\[Am]Nous quittons l'île Saint-Louis pour le para\[G]dis.
\[Am]J'ai trouvé la maison dont nous rêvions \[G]tant,
\[Am]Pour trois fois rien à crédit sur deux \[G]ans.
\endverse
\beginverse
\[Am]C'est au c{\oe}ur du Larzac au bord d'une ri\[G]vière,
\[Am]Dans un joli lieu-dit appelé Le Dé\[G]sert.
\[Am]Pour l'arrivée d'eau le vieux puits fera l'af\[C]faire,
\[G]Pour l'électricité vivent les panneaux so\[Am]laires.
\endverse
\beginverse
\[Am]Le premier spot wifi est à 25 \[G]km,
\[Am]Le premier monop' est à 35 \[G]km,
\[Am]Le premier iPhone est à 120 \[G]km,
\[E]La dernière poste a fer\[Am]mé.
\endverse
\beginchorus
\[Am]Elle est pas belle, la \[C]vie,
\[G]Pour le dernier des hip\[Am]pies ?
\[Am]La main dans la \[C]main
\[G]Avec le dernier la\[Am]pin.
\endchorus
\beginverse
\[Am]Quand on sera vieux on aura tout le \[G]temps
\[Am]De penser au monde qu'on laisse à nos en\[G]fants,
\[Am]Mais là on est trop jeunes et moi j'veux pas cre\[G]ver
\[Am]Trop loin d'un Starbuck ou d'un resto japo\[G]nais.
\endverse
\beginverse
\[Am]Minou fais tes valises et les miennes aus\[C]si,
\[G]Nous quittons le Larzac pour le para\[Am]dis.
\[Am]J'ai trouvé le loft dont tu rêvais \[C]tant,
\[G]Aux pieds de Notre-Dame à crédit sur cent \[Am]ans.
\endverse
\endsong

\beginsong{Le vent des prophètes}[by={Pierre-Michel Gambarelli}]
\beginverse
\[Am]Tu as fait se lever un peuple sans le\[E]vain,
\[Am]Le froment est battu à la \[E]force des \[Am]poings.
\[G]Le pain cuit au soleil a pris le goût du \[C]sang,
\[Dm]Ton nom vient prolonger la liste des gê\[E7]nants.
\endverse
\beginchorus
\[Am]Mais le vent des pro\[G]phètes a souf\[F]flé ce ma\[G]tin, \[E7]
\[Am]Et l'on a vu des milliers d'a\[G]louettes
\[F]Danser \[Em]autour du pè\[Am]lerin.
\[Am]Et l'on a vu sur toute la pla\[G]nète,
\[F]Des \[Em]frères se donner la \[Am]main.
\endchorus
\beginverse
\[Am]L'injustice ne peut supporter le si\[E]lence,
\[Am]Mais l'exemple du Christ prend \[E]toute sa puis\[Am]sance.
\[G]Suspendu au gibet le jour de l'aban\[C]don,
\[Dm]Pour notre humanité il demande par\[E7]don.
\endverse
\beginverse
\[Am]Tu es mort le printemps n'avait que quelques \[E]brins,
\[Am]Tué dans ton église, les \[E]statues pour té\[Am]moins.
\[G]Prêtre d'El Salvador, libre à en mou\[C]rir,
\[Dm]Ton nom vient prolonger la liste des mar\[E7]tyrs.
\endverse
\beginverse
\[Am]Combien d'hommes sont morts et combien meurent en\[E]core,
\[Am]Ils sont pour notre monde cet \[E]envers du dé\[Am]cor.
\[G]Cachés par les parades et les marches mili\[C]taires,
\[Dm]Qui pourrait désormais les forcer à se \[E7]taire ?
\endverse
\endsong
